\documentclass[a4paper,10pt]{article}
\usepackage[cm]{fullpage}
\usepackage{graphicx}
\usepackage{epstopdf}
\usepackage{hyperref}

\begin{document}

\begin{titlepage}
  \setlength{\parindent}{0cm}

  \includegraphics[width=200px]{./Images/uob-logo-black-transparent.eps}

  \Large
  Faculty of Engineering and Design

  \vspace{80pt}

  \LARGE
  Final Year MEng Project \\
  Project Plan and Literature Review

  \vspace{80pt}
  \textbf{Robot Football - High Level Control and Strategy Implementation}

  \vspace{10pt}
  \emph{Anthony Richards} \\
  \emph{\today}

  \vspace{80pt}
  Supervisor: \emph{Dr.~Pejman Iravani}

  \vspace{10pt}
  Assessor: \emph{TBD}
\end{titlepage}

\pagenumbering{roman}

\begin{abstract}
TODO : Abstract
\end{abstract}

\tableofcontents

\cleardoublepage

\pagenumbering{arabic}

\section{Project Background}
TODO : Background

\section{Aims and Objectives}
TODO : Aims and objectives

\section{Methodology and Deliverables}
TODO : Methodology and deliverables

\section{Resources}
As this project is targeted at the SimuroSot league, all the activities will be carried out in software only.  In order to produce the required software, the following software will be required:

\begin{itemize}
 \item Microsoft Visual Studio 2010
 \item SimuroSot Middle League Simulator
\end{itemize}

Visual Studio is available on the University computers in the Maxwell and Faraday labs in the Electrical Engineering Department, as well as on the GigaTerms server which is available via Remote Desktop.  It is also available free of charge to students from the Microsoft Dreamspark Website \cite{dreamsparkVS}. The SimuroSot Simulator is also available free of charge from the FIRA website \cite{simurosotSim}.

In order to run the required software, access to a PC running Microsoft Windows will be required.  These computers are readily available in the Maxwell laboratory in the Electrical Engineering Department, as well as the Software Projects Lab in the same department.  A personal computer capable of running the required software is also available, should it be required.

The development of the project will be aided by the use of an online note-taking website (EverNote), and source code control will be achieved using Subversion.  Combined, these will allow development of the project to be carried on any computer with the suitable software, removing the need for dedicated lab space for the project.  All the required support software is readily available on the University computers and the personal machine.

\section{Predicted Timescale}
TODO : Gannt Chart

\cleardoublepage
\bibliographystyle{plain}
\phantomsection
\addcontentsline{toc}{section}{References}
\bibliography{References/references}

\end{document}
